\subsection*{Chapter 27: Reinforcement Learning with Verifiable Rewards: GRPO}

\paragraph{Exercise 27.1} With partial credit the rewards of \cref{fig:grpo:step} become
$(1,1,0.5,0)$: mean $0.625$, unbiased std $0.479$, advantages
$(+0.783,+0.783,-0.261,-1.305)$. The unboxed ``7'' is still pushed down, but by a third of
its previous weight, and the wrong boxed ``8'' is now pushed down hardest. The loss is
$-\frac14[0.783(-2.269)+0.783(-3.059)-0.261(-1.299)-1.305(-3.400)]=-0.151$. Partial credit
makes \emph{correctness} the dominant signal early (correct-but-unboxed answers are no longer
lumped with wrong ones), which helps when the model is rarely correct. It weakens the
pressure to adopt the format, since an unboxed correct answer is only mildly penalised, so
format adoption is slower. A common compromise is to anneal the partial credit to 0.

\paragraph{Exercise 27.2} With identical rewards, $r-\bar r=0$ for every rollout, the
advantages are exactly 0, the loss is $-0\cdot\sum\log\pi=0$ and autograd propagates zeros to
every parameter (the optimiser step is still taken; AdamW's weight decay is the only change,
and it is zero in our loop). This is desirable because the group carries no information about
\emph{which} behaviour is better. For $p\sim U[0,1]$,
$\E[p^G+(1-p)^G]=\frac2{G+1}$: 40\% wasted for $G=4$ and 11.8\% for $G=16$.

\paragraph{Exercise 27.3} With $k$ correct rollouts, $\bar r=p=k/G$ and
$\sum_i(r_i-\bar r)^2=k(1-p)^2+(G-k)p^2=Gp(1-p)$, so $s^2=\frac{G}{G-1}p(1-p)$, giving
$A^+=(1-p)/s$ and $A^-=-p/s$. Then $\sum_iA_i=[k(1-p)-(G-k)p]/s=G(p-p)/s=0$, and
$\frac1{G-1}\sum_iA_i^2=\frac{Gp(1-p)}{(G-1)s^2}=1$ exactly (ignoring $\epsilon$). As $p\to1/G$
with $G$ large, $A^+=(1-p)/s\approx\sqrt{(1-p)/p}\approx\sqrt{G}$: a single success in a large
group receives an advantage of order $\sqrt G$.

\paragraph{Exercise 27.4} $b_i$ is a function of the other samples only, which are independent
of $y_i$ given $x$, so $\E[b_i\nabla\log\pi(y_i)]=\E[b_i]\,\E[\nabla\log\pi(y_i)]=0$ and the
estimator is unbiased. Algebra: $r_i-b_i=r_i-\frac{G\bar r-r_i}{G-1}=\frac{G}{G-1}(r_i-\bar r)$.
So GRPO's mean-centred advantage equals the RLOO advantage times $\frac{G-1}{G}$, a constant
that the learning rate absorbs. The ``bias'' of including $r_i$ in its own baseline is
therefore only a rescaling. Dividing by the std adds a genuine per-prompt reweighting.

\paragraph{Exercise 27.5} (a) Summed loss: every token of rollout $i$ has gradient weight
$A_i/G$, independent of length. (b) Per-sequence mean: each token has weight
$A_i/(G|y_i|)$. For correct answers of lengths 20 and 2000, the per-token push ratio is
$1:1$ under (a) and $100:1$ under (b). Under (a) the total push on a sequence's log-likelihood
is the same for both, but because $\log\pi$ of the long sequence is a sum of 100 times more
terms, raising it is ``harder''. In practice the summed objective most efficiently increases
the probability of the \emph{short} correct answer. Under (b), long sequences are barely
pushed per token, so long \emph{wrong} answers are barely penalised. That is the length bias
DAPO and Dr.~GRPO criticise (\cref{ch:grpo2}).

\paragraph{Exercise 27.6} PyTorch raises \emph{``Inference tensors cannot be saved for
backward''} when the sampled ids (created inside inference mode) are fed through
\code{nn.Embedding} in the differentiable pass. The embedding saves its index input for
backward. Fixes: sample under \code{no\_grad}; or keep inference mode for sampling but
\code{.clone()} the ids outside it (cloning outside inference mode produces a normal tensor).

\paragraph{Exercise 27.7} $16\times4096\times\num{151936}\times4\,\text{B}\approx39.8$\,GB for
the log-softmax, and the same again for its gradient, so $\approx80$\,GB. Chunked version:
\begin{lstlisting}
def sequence_logprob_chunked(model, ids, prompt_len, C=512):
    h = model_hidden_states(model, ids)          # (T, d) before the LM head
    total = 0.0
    for s in range(prompt_len - 1, ids.numel() - 1, C):
        logits = model.out_head(model.final_norm(h[s:s+C])).float()
        lp = torch.log_softmax(logits, -1)
        total = total + lp.gather(1, ids[s+1:s+1+C, None]).sum()
    return total
\end{lstlisting}
Peak memory scales with $C\cdot|\mathcal V|$ instead of $T\cdot|\mathcal V|$. With
activation checkpointing per chunk the backward is bounded the same way. (Exposing hidden
states requires a small model hook; \code{llmsys.rl.logprobs.token\_logprobs(chunk=...)} chunks
after the LM head, which bounds the float32 copy but not the bf16 logits.)

\paragraph{Exercise 27.8} Sequentially, $t_{\text{sample}}$ grows linearly in both $G$ and $T$
(one forward per token per rollout); on CPU the toy gives $a$ of the order of milliseconds per
token. $t_{\text{train}}$ is also linear in $G\cdot T$ but with a smaller constant per token,
since a prefill-style forward processes all tokens in parallel. With \code{sample\_rollouts}
the $G$ rollouts share each decode step, so $t_{\text{sample}}\approx a'\,T$ with $a'$ almost
independent of $G$ until the batch stops being memory-bound. This is the decode roofline:
weights are streamed once per step regardless of batch size.

\paragraph{Exercise 27.9} Qwen3-4B: 36 layers, 8 KV heads, head dim 128, so the KV cache is
$2\cdot36\cdot8\cdot128\cdot2=147$\,KB per token. (i) $256$ sequences $\times4352$ tokens
$\approx1.11\times10^6$ tokens, about 164\,GB. (ii) Weights $\approx8.0$\,GB in bf16. Each of the
4096 decode steps streams the weights plus a cache that grows from 9.7\,GB to 164\,GB, about
95\,GB on average, so roughly $4096\times95\,\text{GB}/3.35\,\text{TB/s}\approx116$\,s.
(iii) Training FLOPs $6\cdot4.0\times10^9\cdot1.11\times10^6\approx2.7\times10^{16}$, about 68\,s at
40\% MFU. The cache does not fit next to the weights on one 80\,GB GPU ($\approx72$\,GB free).
It needs $\lceil164/72\rceil=3$ waves (or 3+ GPUs), which roughly preserves total sampling
time, since each wave streams the weights 4096 times.

\paragraph{Exercise 27.10} Examples: (1) ``\verb|\boxed{7}| \ldots\ actually \verb|\boxed{8}|''
(first-box readers accept the wrong final answer); (2) ``\verb|\boxed{17}|'' with truth 7
(substring match accepts); (3) ``\verb|\boxed{7.0}|'' (no normalisation rejects a correct
answer); (4) ``\verb|\boxed{\frac{14}{2}}|'' (rejected without fraction handling);
(5) ``\verb|\boxed{1, 2, 3, 4, 5, 6, 7}|'' (substring or ``contains'' accepts). Harden by
taking the last box, comparing normalised values exactly, and rejecting boxed lists when a
scalar is expected. RL is most likely to discover the cheap ones (many boxes, lists), because
they increase reward on \emph{every} prompt without solving anything.

\paragraph{Exercise 27.11} With 20 held-out facts, GRPO transfers the \emph{format} almost
perfectly (boxing is prompt-independent behaviour). Correctness on held-out facts improves
little, because a 76k-parameter model trained on 80 facts has not learned an addition
algorithm, and RL with an outcome reward cannot teach a fact it never produces. This mirrors
the MATH/MATH-500 split: RLVR mostly improves how existing knowledge is deployed (format,
persistence, self-checking), and held-out evaluation is what separates that from memorising
the training prompts.
